\section{Limitations and Ethical Considerations}
\label{sec:limitations}

The live evaluation uses \V{lg.clips} clips recorded by a single signer, the author, in one room
with one laptop camera. The author is not a fluent signer; the signs in the live recordings were
learned from online ASL video dictionaries, and how the \V{lg.signs} signs were chosen was not
recorded. If the signer's productions differ from the dataset's signers, some live errors are
production differences rather than pipeline or camera effects, and nothing in this design
separates them. The results show which pipeline discrepancies matter for this signer and camera;
they do not estimate accuracy for Deaf signers in general. Two controls would locate the remaining
gap: running the live pipeline on test-split landmark sequences, and an oracle segmentation of the
live clips. The ASL Citizen label mapping is AI-assisted and unconfirmed by a signer, and v3's
in-domain accuracy does not validate it. The live clips are not an independent test of v3 either:
they were analyzed before v3's webcam augmentation was designed. The model recognizes
\V{data.signs} isolated signs; it does not handle continuous signing, fingerspelling, or signs
outside the vocabulary, and it is not a translator. The author is based in India, where Indian
Sign Language is the relevant language; nothing here transfers to it without new data.

We did not evaluate the system with Deaf or hard-of-hearing users, and no Deaf signer took part in
its design. Sign language technology should be built with, and ideally led by, Deaf communities
\cite{bragg2019interdisciplinary,yin2021including,fox2023bestpractices}, and Deaf researchers have
criticized how often the field sets its agenda without them \cite{desai2024systemic}. At the live
accuracy we measured, the app's Practice mode, which scores a sign against a target, would tell a
learner that many correctly produced signs are wrong; it should be evaluated with learners and
Deaf teachers before anyone relies on it. The live clips show the author, who consents to their use here; they are not
released, and no ASL Citizen data or model trained on it is released either. When the user asks
the app for a sentence, only the recognized glosses are sent to a third-party language model
provider. The code, configurations and result files behind every number are in the repository of
the extended version \cite{shukla2026manuvox}.
